\textcolor{orange}{This is a placeholder section, I'm writing down a description for each attribute and each data set that i used. I will then need to structure it properly to make it readable.}

\subsubsection{Thermal Efficiency}
Thermal efficiency indicates the fraction of thermal power generated by fission that is converted into net electrical power, accounting for in-plant consumption. Efficiency values were calculated using net electrical power output and thermal power data from the ARIS database \cite{iaea2025aris}:
$$
\eta = \frac{P_{el,net}}{P_{th}}
$$

\begin{table}[!ht]
    \centering
    \begin{tabularx}{\textwidth}{X c}
        \toprule
        \textbf{Reactor Model} & \textbf{Net Efficiency [\%]} \\
        \midrule
        APR1400         & 35.2 \\
        EPR 1750        & 36.2 \\
        AP1000          & 32.4 \\
        BWRX-300        & 34.5 \\
        Rolls Royce SMR & 34.6 \\
        NuScale         & 30.8 \\
        \bottomrule
    \end{tabularx}
    \caption{Net electrical efficiencies for selected reactor models}
    \label{tab:reactor-efficiency}
\end{table}



\subsubsection{Discharge Burnup}
Discharge burnup measures the energy extracted per unit of heavy metal fuel ($MWd/tonHM$). Data for APR1400 and AP1000 were sourced from the ARIS database \cite{iaea2025aris}, while EPR 1750 data were obtained from \cite{NRC_ML052280170}. Ranges are provided for SMRs due to design variability.

\begin{table}[!ht]
    \centering
    \begin{tabularx}{\textwidth}{X c}
        \toprule
        \textbf{Reactor Model} & \textbf{Discharge Burnup [$MWd/tonHM$]} \\
        \midrule
        APR1400         & 46.5 \\
        EPR 1750        & 55 \\
        AP1000          & 50 \\
        BWRX-300        & $50 \pm 5$ \\
        Rolls Royce SMR & $55 \pm 5$ \\
        NuScale         & $52.5 \pm 2.5$ \\
        \bottomrule
    \end{tabularx}
    \caption{Discharge burnup for selected reactor models}
    \label{tab:reactor-burnup}
\end{table}





\subsubsection{Technological Readiness and Planned Commissioning Time}
The technological Readiness Level is a purely qualitative score that is given on a scale from 1 to 13 as discrete values from the Fibonacci sequence (1,2,3,5,8,13), where 1 is the furtherst away from commercial deplyment and 13 is the closest to Nth-of-a-kind. The score considers the state of development, current number of units operational or under construction and licensing status as well as other considerations. The data was gathered from the ARIS database \cite{iaea2025aris} and from the respective company websites. \\

The APR1400 and the AP1000 are considered the most advanced reactors, with several units operational (9 and 7 respectively) and other under construction (3 and 2). The only catch with these designs is that none of these deployments has been made within Europe. On the other hand the EPR 1750 has only 2 operational units, both in China, 2 similar operational units (EPR 1650 in Olkiluoto and Flamanville) and 3 units Under construction in the UK (Hinkley Point and Sizewell). So these three reactors all get a 8/13 score, the first and the second don't get full points because of the lack of European experience, while the EPR 1750 doesn't get full points because of the limited number of operational units and the delays and overcosts that have plagued the first two EPRs. \\

Among the three SMRs the Rolls Royce SMR is the most similar to a conventional power plant being based on a multi-loop design with forced circulation, the know-how trasnferability for commissioing is considered high, as well as the managment of the construction site. For this it scores a 5/13. \\

The BWRX-300 is a more innovative design being an integrated reactor housing all the components of the NSSS within the reactor vessel, this means that the construction itself will be the first of his kind but this type of desing has been studied since 2010s \cite{PLACEHOLDER}. Moreover since the plan is to deploy just one unit the construction site will most likely be managed as a common nuclear power plant. For tthis it scores a 3/13. \\

At last, the Nuscale SMR has a compleltly novel design, not only is a natural circulation based, integrated PWR, but it is also ment to be deployed in a multi-unit configuration in a novel way. This means that the managment of the construction site itself will be new, for this it scores the lowest, a 2/13. \\



\begin{table}[!ht]
    \centering
    \begin{tabularx}{\textwidth}{l X X}
        \toprule
        \textbf{Reactor Model} & \textbf{Net Power Output $[MW_{e,net}]$} & Data Source\\
        \midrule
        APR1400         & 1340 & \cite{IAEA_PRIS_2025} \\
        EPR 1750        & 1630 & \cite{IAEA_PRIS_2025} \\
        AP1000          & 1120 & \cite{IAEA_PRIS_2025} \\
        BWRX-300        & 300 & \cite{iaea2025aris} \\
        Rolls Royce SMR & 470 & \cite{iaea2025aris} \\
        NuScale         & 77 (per unit) & \cite{iaea2025aris} \\
        \bottomrule
    \end{tabularx}
    \caption{Load following capabilities and power output}
    \label{tab:load-power_output}
\end{table}



\subsubsection{Radioactive Waste Produced}
TD

% Results are shown in Table \ref{tab:waste-produced}.
% \begin{table}[!ht]
%     \centering
%     \begin{tabularx}{0.6\textwidth}{l c}
%         \toprule
%         \textbf{Reactor Model} & \textbf{Waste Produced [kg/TWh]} \\
%         \midrule
%         APR1400         & 22 \\
%         EPR 1750        & 20 \\
%         AP1000          & 24 \\
%         BWRX-300        & 28 \\
%         Rolls Royce SMR & 30 \\
%         NuScale         & 35 \\
%         \bottomrule
%     \end{tabularx}
%     \caption{Radioactive waste produced per TWh}
%     \label{tab:waste-produced}
% \end{table}



\subsection{Country-Specific and Contextual Factors}
This subsection evaluates factors dependent on the deployment country, including fuel availability, regulatory environment, and supply chain adaptability.

\subsubsection{Fuel Manufacturing Feasibility by Country}
For this we have to consider two aspects, the first one is technolgy-wise and considers if the fuel exists or not. The second cosnideration is dependent on the coutry and considers if the that fuel is used in that country or how easy would it be to get to that country.

For the country aspect we can apply a correction factor as follows: 
$1$ is the fuel is already used in that country such as EPR fuel in France. $0.9$ if the fuel is used in other european countries since some paperwork will be needed but the open market will facilitate the import. $0.8$ for the same situation but for CH since the market has some barriers. $0.5$ is that fuel is not at all used in Europe and has to be imported from overseas.

Then regarding the tecnology, the 3 existing reactors get a 10/10 score, as well as the BWRX-300 since it will use an existing fuel assembly GNF-2 \cite{CNSC_BWRX300_Fuel_Design}. Rolls-Royce SMR and NuScale get a 5/10 since they will have to develop their own fuel but it will be a derivative of existing PWR fuel \cite{RollsRoyce_SMR_UK_Fuel_2023} and \cite{NRC_ML17007A001}. Rolls-Royce is working with Westinghouse for the fuel development \cite{PLACEHOLDER} while NuScale wants to modify the HTP\texttrademark{} fuel that is also produced by Framatome \cite{Framatome_GNF2_Fuel_2018}.

% So here there is a cross table with the reactors and the countries with the respective scores.

% \begin{table}[!ht]
%     \centering
%     \begin{tabularx}{0.9\textwidth}{l c c c c c}
%         \toprule
%         \textbf{Reactor Model} & \textbf{France (FR)} & \textbf{Switzerland (CH)} & \textbf{Poland (PO)} & \textbf{Northern Italy (NI)} & \textbf{Base Score} \\
%         \midrule
%         EPR1750         & 1.0 & 1.0 & 1.0 & 1.0 & 10 \\
%         APR1400         & 1.0 & 1.0 & 1.0 & 1.0 & 10 \\
%         AP1000          & 1.0 & 1.0 & 1.0 & 1.0 & 10 \\
%         BWRX-300        & 1.0 & 1.0 & 1.0 & 1.0 & 10 \\
%         Rolls Royce SMR & 1.0 & 1.0 & 1.0 & 1.0 & 5 \\
%         NuScale         & 1.0 & 1.0 & 1.0 & 1.0 & 5 \\
%         \bottomrule
%     \end{tabularx}
%     \caption{Fuel manufacturing feasibility scores by country}
%     \label{tab:fuel-feasibility}
% \end{table}



\subsubsection{Know-how Transferability and Supply Chain Adaptability}
Scores reflect the ease of transferring operational knowledge and adapting local supply chains. Know-how transferability considers the presence of similar designs in the country whereas supply chain adaptability considers the presence of the NSSS supplier in the country. For instance France has a lot of experience with PWR and has an EPR (of previous iteration at Flamanville) therefore will have a high know-how transferability for the EPR, on the other hand Italy has shut down all its reactors since 1986 therefore the know-how related score will be low for all reactors. Nonetheless Italy still has an active Nuclear Industry, with Ansaldo Nucleare operating in Genova and Framatome having offices in Milan. This means that even if the know-how is low, the supply chain can be esablished more easily. \\

The reason to separate these two aspects is to be able to give different importance (weights) to the human aspect (know-how) and the industrial/Economic aspect of the supply chain. \\

% The results are shown in Table \ref{tab:know-how-country} and Table \ref{tab:supply-chain-country}.

% \begin{table}[!ht]
%     \centering
%     \begin{minipage}{0.48\textwidth}
%         \centering
%         \begin{tabularx}{\textwidth}{l c c c c}
%             \toprule
%             \textbf{Reactor Model} & \textbf{IT} & \textbf{FR} & \textbf{CH} & \textbf{PO} \\
%             \midrule
%             APR1400         & 4 & 8 & 7 & 6 \\
%             EPR 1750        & 5 & 10 & 8 & 7 \\
%             AP1000          & 4 & 7 & 7 & 6 \\
%             BWRX-300        & 3 & 6 & 6 & 5 \\
%             Rolls Royce SMR & 3 & 6 & 5 & 5 \\
%             NuScale         & 2 & 5 & 4 & 4 \\
%             \bottomrule
%         \end{tabularx}
%         \caption{Know-how transferability scores by country}
%         \label{tab:know-how-country}
%     \end{minipage}
%     \hfill
%     \begin{minipage}{0.48\textwidth}
%         \centering
%         \begin{tabularx}{\textwidth}{l c c c c}
%             \toprule
%             \textbf{Reactor Model} & \textbf{IT} & \textbf{FR} & \textbf{CH} & \textbf{PO} \\
%             \midrule
%             APR1400         & 7 & 8 & 7 & 7 \\
%             EPR 1750        & 8 & 9 & 8 & 8 \\
%             AP1000          & 7 & 7 & 7 & 7 \\
%             BWRX-300        & 6 & 6 & 6 & 6 \\
%             Rolls Royce SMR & 6 & 7 & 6 & 6 \\
%             NuScale         & 5 & 6 & 5 & 5 \\
%             \bottomrule
%         \end{tabularx}
%         \caption{Supply chain adaptability scores by country}
%         \label{tab:supply-chain-country}
%     \end{minipage}
% \end{table}



\subsection{Economic and Political Indicators}
This subsection assesses macro-level indicators influencing project feasibility.


\subsubsection{Marginal Cost Impact}
TD

% Table \ref{tab:marginal-cost-impact} summarizes the qualitative impact.
% \begin{table}[!ht]
%     \centering
%     \begin{tabularx}{0.5\textwidth}{l c}
%         \toprule
%         \textbf{Reactor Model} & \textbf{Marginal Cost Impact} \\
%         \midrule
%         APR1400         & Low \\
%         EPR 1750        & Moderate \\
%         AP1000          & Low \\
%         BWRX-300        & High \\
%         Rolls Royce SMR & High \\
%         NuScale         & High \\
%         \bottomrule
%     \end{tabularx}
%     \caption{Marginal cost impact for selected reactor models}
%     \label{tab:marginal-cost-impact}
% \end{table}



\subsubsection{Country Economics and Politics}
These indicators are meant to represent a country's economic and political state. 
Real GDP from 2015 to 2030 (projected) and Expense as a percentage of GDP from 2015 to 2023 have been obtained from the International Monetary Fund \cite{PLACEHOLDER}. \\

Researchers in Research and Developments per million inhabitants from 2015 to 2023, as well as the GDP expenditure for Research have been obtained from the UNESCO Institute for Statistics \cite{PLACEHOLDER}. \\
Political Stability, a qualitative indicator ranging from $-2.5$ to $2.5$, has been obtained from the Worldwide Governance Indicators project by the World Bank \cite{PLACEHOLDER}. \\
Then from the V-Dem dataset \cite{PLACEHOLDER} we extracted:
\begin{itemize}
    \item v2x polyarchy: TO what extent is the people's opinion taken into account in the political process?
    \item v2regdur: Regime duration, how long has the last regime (goverment) lasted?
    \item v2cltrnslw: Are the laws and the political actions transparent, well publicized, coherent and stable thoughout the years?
    \item v2dlengage: To what extent are the citizens engaged in the political process?
\end{itemize}

Then we considered public support for nuclear energy, data was gathered from various sources \cite{PLACEHOLDER} \cite{PLACEHOLDER} \cite{PLACEHOLDER} \cite{PLACEHOLDER}.

The objective of these data is to create 2 contry-specific modifiers. One for their economic state, if a country is growing and tends to spend a lot in R\&D it will be more likely to be able to afford a nuclear project. The second indicator is to represent the support for nuclear energy, not only considering the current stance, but also trying to understand how stable this stance will be in the coming years.

% \begin{table}[!ht]
%     \centering
%     \begin{tabularx}{0.8\textwidth}{l c c c}
%         \toprule
%         \textbf{Country/Region} & \textbf{Real GDP Growth [\%]} & \textbf{Political Stability Score} & \textbf{Public Support for Nuclear [\%]} \\
%         \midrule
%         France          & 1.8 & 0.9 & 65 \\
%         Switzerland     & 1.2 & 0.8 & 45 \\
%         Poland          & 3.5 & 0.7 & 70 \\
%         Northern Italy  & 0.9 & 0.6 & 35 \\
%         \bottomrule
%     \end{tabularx}
%     \caption{Country economics and politics indicators}
%     \label{tab:country-indicators}
% \end{table}

\subsubsection{Electric data}
TD

\subsubsection{Environmental Impact: GHG Emission Intensity Benefit}

\subsubsection{GHG Emission Reduction Estimation}
Greenhouse gas (GHG) emission reductions were quantified by comparing each country's 2024 emissions against projected emissions following deployment of one reactor unit. Country-specific GHG emission intensities were derived from Electricity Maps \cite{ElectricityMaps_CarbonIntensity_2024}. Total load and cross-border flow measurements were obtained from the ENTSO-E Transparency Platform \cite{ENTSOE_Transparency_2025}. Due to non-normal data distributions, median values with interquartile ranges were employed to characterize uncertainty (Table~\ref{tab:ghg-intensity-country}). \\

For nuclear generation's lifecycle emissions, Liu et al.'s (2023) comprehensive meta-analysis of 42 studies \cite{10.3389/fenrg.2023.1147016} served as the primary reference, reporting a global average of $15.815 \pm 3.59 \; gCO_{2,\text{eq}}/kWh_e$ for light water reactors. This conservative estimate was selected despite lower values observed in specific contexts: French nuclear operations typically range between $ 2 \sim 6 \; gCO_{2,\text{eq}}/kWh_e$ according to both operational data \cite{paillere_2021} and IAEA assessments \cite{IAEA_Climate_Change_Nuclear_Power_2022}, while individual studies on french reactors have reported values as high as $12 \; gCO_{2,\text{eq}}/kWh_e$ \cite{paillere_2021}. The global average was preferred to account for potential regional variations in fuel cycle practices and construction impacts. \\

Small modular reactors (SMRs) were assigned an emission intensity of $17.24 \pm 3.91$ $gCO_{2,\text{eq}}/kWh_e$, incorporating Carless et al.'s (2016) finding that SMRs exhibit approximately $9\%$ higher lifecycle emissions than conventional large LWRs due to reduced economies of scale in construction and fuel production \cite{CARLESS201684}. \\

The GHG reduction benefit was calculated as the percentage improvement between current energy mixes $GHG_{current\ mix}$  and the nuclear deployment scenario $GHG_{nuclear}$:
$$
GHG_{Benefit} (\%) = \left(1 - \frac{GHG_{nuclear}}{GHG_{current\ mix}}\right) \times 100
$$

% The results are shown in Table \ref{tab:ghg-benefit-country}.
% \begin{table}[!ht]
%     \centering
%     \begin{tabularx}{0.7\textwidth}{l c c}
%         \toprule
%         \textbf{Country/Region} & \textbf{GHG Reduction (Large) [tCO2/MWh]} & \textbf{GHG Reduction (SMR) [tCO2/MWh]} \\
%         \midrule
%         Italy (IT)        & 0.85 & 0.83 \\
%         Switzerland (CH)  & 0.10 & 0.09 \\
%         Poland (PO)       & 0.90 & 0.88 \\
%         France (FR)       & 0.15 & 0.14 \\
%         \bottomrule
%     \end{tabularx}
%     \caption{GHG emission intensity benefit by country for large reactors and SMRs}
%     \label{tab:ghg-benefit-country}
% \end{table}



\subsubsection{Import Dependence}
TD

\subsubsection{Policy Stability and Long-Term Project Viability}





